\subsection{Haar measure of an inverse limit*}

Let G be a locally compact group (hence complete). Let \((\mathrm{K}_{\alpha})_{\alpha\in\mathrm{A}}\) be a decreasing directed family of compact normal subgroups of G, with intersection \(\{e\}\) (so that the filter base formed by the \(K_{\alpha}\) converges to e). Set \(G_{\alpha}=G/K_{\alpha}\) ; let \(\varphi_{\alpha}:G\to G_{\alpha}\) and \(\varphi_{\beta\alpha}:G_{\alpha}\to G_{\beta}\) ( \(\alpha\geqslant\beta\) ) be the canonical homomorphisms. Then, the inverse limit of the inverse system \((\mathrm{G}_{\alpha},\varphi_{\beta\alpha})\) may be identified with G, and the canonical mapping of this inverse limit into \(G_{\alpha}\) is identified with \(\varphi_{\alpha}\) (GT, III, §7, No.~3, Prop. 2). The mappings \(\varphi_{\alpha}\) and \(\varphi_{\beta\alpha}\) are proper (loc. cit., §4, No.~1, Cor. 2 of Prop. 1). These assumptions remain fixed throughout this subsection.

Lemma 2. --- a) Let \(f \in \mathcal{K}_{+}(\mathrm{G})\) , S a compact subset of G containing Supp \(f\) , U an open neighborhood of S in G, and \(\varepsilon > 0\) . There exist an \(\alpha \in \mathrm{A}\) and a function \(g \in \mathcal{K}_{+}(\mathrm{G})\) , zero outside U and constant on the cosets of \(\mathrm{K}_{\alpha}\) , such that \(|f - g| \leqslant \varepsilon\) .

\begin{enumerate}
\def\labelenumi{\alph{enumi})}
\setcounter{enumi}{1}
\tightlist
\item
  Let \(\mu\) and \(\mu'\) be two measures on \(\mathbf{G}\) such that \(\varphi_{\alpha}(\mu) = \varphi_{\alpha}(\mu')\) for all \(\alpha \in \mathbf{A}\) . Then \(\mu = \mu'\) .
\end{enumerate}

There exists an \(\alpha_{1} \in \mathbf{A}\) such that \(\mathrm{K}_{\alpha_1}\mathrm{S} \cap \mathrm{K}_{\alpha_1}(\mathrm{G} - \mathrm{U}) = \varnothing\) (GT, II, §4, No.~3, Prop. 4). Augmenting S and diminishing U, we may therefore assume that S and U are unions of cosets of \(\mathrm{K}_{\alpha_1}\) . Consider the continuous numerical functions \(h\) on S having the following property: there exists an \(\alpha \geqslant \alpha_{1}\) such that \(h\) is constant on the cosets of \(\mathrm{K}_{\alpha}\) . These functions form a subalgebra of \(\mathcal{K}(\mathrm{S})\) (because \((\mathrm{K}_{\alpha})\) is a decreasing directed family) that contains the constants and separates the points of S: for, let \(x, y\) be two distinct points of S; since the intersection of the \(\mathrm{K}_{\alpha}\) is \(\{e\}\) , there exists an \(\alpha \geqslant \alpha_{1}\) such that \(\varphi_{\alpha}(x) \neq \varphi_{\alpha}(y)\) , then a numerical function \(u\) continuous on \(\varphi_{\alpha}(\mathrm{S})\) such that \(u(\varphi_{\alpha}(x)) \neq u(\varphi_{\alpha}(y))\) . By the Stone-Weierstrass theorem, there exist an \(\alpha \geqslant \alpha_{1}\) and a continuous function \(h \geqslant 0\) on S, constant on the cosets of \(\mathrm{K}_{\alpha}\) , such that \(|f - h| \leqslant \frac{\varepsilon}{2}\) on S. For every \(t \in \mathbf{R}\) , set \(\delta(t) = \left(t - \frac{\varepsilon}{2}\right)^+\) , and set \(h' = \delta \circ h\) . Then \(h'\) is a function \(\geqslant 0\) , continuous on S, constant on the cosets of \(\mathrm{K}_{\alpha}\) , and \(|h - h'| \leqslant \frac{\varepsilon}{2}\) on S, therefore \(|f - h'| \leqslant \varepsilon\) on S. On the other hand, \(h'(x) = 0\) if \(x\) belongs to the boundary of S in G, because then \(h(x) \leqslant \frac{\varepsilon}{2}\) . If \(h'\) is extended by 0 on the complement of S, one obtains a function \(g\) that meets the requirements, which proves a).

Now let \(\mu, \mu'\) be two measures on \(G\) such that \(\varphi_{\alpha}(\mu) = \varphi_{\alpha}(\mu')\) for all \(\alpha \in A\) . Let \(v \in \mathcal{K}(G)\) be a function constant on the cosets of \(K_{\alpha}\) for some \(\alpha \in A\) , so that we may write \(v = w \circ \varphi_{\alpha}\) with \(w \in \mathcal{K}(G_{\alpha})\) ; then \(\mu(v) = (\varphi_{\alpha}(\mu))(w) = (\varphi_{\alpha}(\mu'))(w) = \mu'(v)\) ; it follows that \(\mu = \mu'\) by virtue of a).

PROPOSITION 6. --- For every \(\alpha \in \mathbf{A}\) , let \(\mu_{\alpha}\) be a positive measure on \(\mathbf{G}_{\alpha}\) . Suppose that \(\varphi_{\beta \alpha}(\mu_{\alpha}) = \mu_{\beta}\) for \(\alpha \geqslant \beta\) . Then, there exists one and only one positive measure \(\mu\) on \(\mathbf{G}\) such that \(\varphi_{\alpha}(\mu) = \mu_{\alpha}\) for all \(\alpha \in \mathbf{A}\) .

Uniqueness follows immediately from Lemma 2 b). Let us prove the existence of \(\mu\) . Let V be the vector space of functions belonging to \(\mathcal{K}(G)\) and constant on the cosets of some \(K_{\alpha}\) ( \(\alpha\) may depend on the function). It follows from Lemma 2 a) that V satisfies the condition (P) of Ch. III, §1, No.~7, Prop. 9: for, let K be a compact set in G and choose \(f \in \mathcal{K}_{+}(G)\) with \(f(x) > 0\) for all \(x \in K\) ; let \(a > 0\) be the smallest value of \(f\) on K; by Lemma 2 a), there exists a function \(g \in V \cap \mathcal{K}_{+}(G)\) such that \(|f - g| \leqslant a/2\) , therefore \(g(x) > 0\) for all \(x \in K\) , and condition (P) is verified. Let \(f \in V\) . There exists an \(\alpha \in A\) such that \(f\) is constant on the cosets of \(K_{\alpha}\) . By passage to the quotient, \(f\) defines a function \(f_{\alpha} \in \mathcal{K}(G_{\alpha})\) . The number \(\mu(f) = \mu_{\alpha}(f_{\alpha})\) does not depend on the choice of \(\alpha\) : for, let \(\beta\) be any index such that \(f\) is constant on the cosets of \(K_{\beta}\) ; let \(\gamma \in A\) be such that \(\gamma \geqslant \alpha\) , \(\gamma \geqslant \beta\) ; then \(f\) defines functions \(f_{\beta} \in \mathcal{K}(G_{\beta})\) , \(f_{\gamma} \in \mathcal{K}(G_{\gamma})\) such that \(f = f_{\beta} \circ \varphi_{\beta} = f_{\gamma} \circ \varphi_{\gamma}\) ; then \(f_{\alpha} \circ \varphi_{\alpha\gamma} = f_{\gamma}\) , therefore \(\mu_{\gamma}(f_{\gamma}) = (\varphi_{\alpha\gamma}(\mu_{\gamma}))(f_{\alpha}) = \mu_{\alpha}(f_{\alpha})\) , and similarly \(\mu_{\gamma}(f_{\gamma}) = \mu_{\beta}(f_{\beta})\) , whence our assertion. This established, it is clear that \(\mu\) is a linear form on V and that \(\mu(f) \geqslant 0\) for \(f \geqslant 0\) . By Prop. 9 of Ch. III, §1, No.~7, \(\mu\) may be extended to a positive measure on G, which we again denote by \(\mu\) . One has \(\varphi_{\alpha}(\mu) = \mu_{\alpha}\) for all \(\alpha \in A\) by the very construction of \(\mu\) .

DEFINITION 5. --- The measure \(\mu\) is said to be the inverse limit (or projective limit) of the \(\mu_{\alpha}\) .

PROPOSITION 7. --- We retain the notations of Prop. 6. If each \(\mu_{\alpha}\) is a left (resp. right) Haar measure on \(G_{\alpha}\) , then \(\mu\) is a left (resp. right) Haar measure on \(G\) .

Suppose, for example, that the \(\mu_{\alpha}\) are left Haar measures. Let \(s \in \mathbf{G}\) . For every \(x \in \mathbf{G}\) ,

\[
\left(\varphi_ {\alpha} \circ \boldsymbol {\gamma} (s)\right) (x) = \varphi_ {\alpha} (s x) = \varphi_ {\alpha} (s) \varphi_ {\alpha} (x) = \left(\boldsymbol {\gamma} (\varphi_ {\alpha} (s)) \circ \varphi_ {\alpha}\right) (x);
\]

therefore \(\varphi_{\alpha}(\pmb{\gamma}(s)\mu) = \pmb{\gamma}(\varphi_{\alpha}(s))\mu_{\alpha} = \mu_{\alpha}\) . Therefore \(\pmb{\gamma}(s)\mu = \mu\) by Lemma 2 b), thus \(\mu\) is a left Haar measure.

We suppose henceforth the \(K_{\alpha}\) to be not only compact, but open in G. The \(G_{\alpha}\) are then discrete and, for \(\beta \geqslant \alpha\) , \(K_{\alpha}/K_{\beta}\) is a compact and discrete group, hence is finite. The group G is unimodular (Prop. 3).

PROPOSITION 8. --- a) Let \(\mu\) and \(\mu'\) be two positive measures on G such that, for every \(\alpha\) and for every coset C of \(\mathrm{K}_{\alpha}\) , \(\mu(\mathrm{C}) = \mu'(\mathrm{C})\) . Then \(\mu = \mu'\) .

\begin{enumerate}
\def\labelenumi{\alph{enumi})}
\setcounter{enumi}{1}
\tightlist
\item
  Fix an \(\alpha_0 \in \mathbf{A}\) . For every \(\alpha \geqslant \alpha_0\) let \(n_{\alpha}\) be the number of elements of the finite group \(\mathrm{K}_{\alpha_0} / \mathrm{K}_{\alpha}\) . There exists one and only one positive measure \(\mu\) on \(\mathbf{G}\) such that, for every \(\alpha \geqslant \alpha_0\) , each coset of \(\mathrm{K}_{\alpha}\) has measure \(n_{\alpha}^{-1}\) . Moreover, \(\mu\) is a Haar measure on \(\mathbf{G}\) , such that \(\mu(\mathrm{K}_{\alpha_0}) = 1\) .
\end{enumerate}

Let \(\mu\) and \(\mu'\) be two positive measures on G satisfying the condition a). The points of the discrete group \(G_{\alpha}\) then have the same measure for \(\varphi_{\alpha}(\mu)\) and \(\varphi_{\alpha}(\mu')\) , whence \(\varphi_{\alpha}(\mu) = \varphi_{\alpha}(\mu')\) , and this for all \(\alpha\) . Therefore \(\mu = \mu'\) (Lemma 2 b)).

Let us prove b). For every \(\alpha \geqslant \alpha_0\) , let \(\mu_{\alpha}\) be the Haar measure of the discrete group \(G_{\alpha}\) such that every point has measure \(n_{\alpha}^{-1}\) . Let \(\alpha, \beta\) be such that \(\alpha \geqslant \beta \geqslant \alpha_0\) . Then \(K_{\beta}/K_{\alpha}\) has \(n_{\alpha}/n_{\beta}\) elements. Therefore \(\varphi_{\beta\alpha}(\mu_{\alpha})\) is the measure on \(G_{\beta}\) such that each point has measure \(n_{\alpha}^{-1} \cdot \frac{n_{\alpha}}{n_{\beta}} = n_{\beta}^{-1}\) ; in other words, \(\varphi_{\beta\alpha}(\mu_{\alpha}) = \mu_{\beta}\) . It then suffices to apply Props. 6 and 7.

Example. --- Let \(Q_{p}\) be the p-adic field, the completion of Q for the p-adic absolute value \(|x|_{p} = p^{-v_{p}(x)}\) (GT, IX, §3, No.~2, Example 3). The elements of \(Q_{p}\) are called p-adic numbers. We denote again by \(|x|_{p}\) the continuous extension of the p-adic absolute value to \(Q_{p}\) . One has

\[
| x + y | _ {p} \leqslant \sup (| x | _ {p}, | y | _ {p})
\]

for \(x, y\) in \(\mathbf{Q}\) (loc. cit.), hence for \(x, y\) in \(\mathbf{Q}_p\) ; moreover, if \(|y|_p < |x|_p\) then \(|x + y|_p = |x|_p\) , because \(|x|_p = |(x + y) - y|_p \leqslant \sup(|x + y|_p, |y|_p)\) . If \((x_n)\) is a sequence of points of \(\mathbf{Q}_p\) tending to \(x \in \mathbf{Q}_p^*\) , then \(|x - x_n|_p < |x|_p\) and \(|x - x_n|_p < |x_n|_p\) for \(n\) sufficiently large, therefore \(|x|_p = |x_n|_p\) . This proves that, for every \(x \in \mathbf{Q}_p^*\) , \(|x|_p\) is a power of \(p\) .

Let \(\mathbf{Z}_p\) be the closure of \(\mathbf{Z}\) in \(\mathbf{Q}_p\) ; this is a subring of \(\mathbf{Q}_p\) ; its elements are called \(p\) -adic integers. One has \(|x|_p \leqslant 1\) for every \(x \in \mathbf{Z}_p\) . Conversely, let \(x\) be an element of \(\mathbf{Q}_p\) such that \(|x|_p \leqslant 1\) , and let us show that \(x \in \mathbf{Z}_p\) ; there exists a sequence \((x_n)\) of elements of \(\mathbf{Q}\) tending to \(x\) , and \(|x_n|_p \leqslant 1\) for \(n\) sufficiently large by what we have seen above; it suffices to show that \(x_n\) belongs to \(\mathbf{Z}_p\) for \(n\) sufficiently large; in other words, we are reduced to the case that \(x \in \mathbf{Q}\) ; then \(x = a/b\) with \(b\) relatively prime to \(p\) ; for every integer \(n > 0\) , there exist \(b_n' \in \mathbf{Z}\) and \(h_n \in \mathbf{Z}\) such that \(bb_n' + h_np^n = 1\) , whence \(x = \frac{abb_n' + ah_np^n}{b} = ab_n' + \frac{ah_np^n}{b}\) and \(|x - ab_n'|_p \leqslant p^{-n}\) , therefore \(ab_n'\) tends to \(x\) .

From this, it follows that the closed ball with center 0 and radius \(p^{-n}\) , identical to the open ball with center 0 and radius \(p^{-n+1}\) , is \(p^n\mathbf{Z}_p\) . The topological space \(\mathbf{Q}_p\) is therefore zero-dimensional, hence totally disconnected (GT, IX, §6, No.~4).

Let us show that the integers \(0,1,\ldots,p^{n}-1\) constitute a system of representatives of \(\mathbf{Z}_{p}\) modulo \(p^{n}\mathbf{Z}_{p}\) . First, \(|k-k^{\prime}|_{p}>p^{-n}\) for two such integers \(k\) and \(k^{\prime}\) , therefore the classes modulo \(p^{n}\mathbf{Z}_{p}\) of these integers are distinct. On the other hand, let \(x\in\mathbf{Z}_{p}\) ; there exists a \(k\in\mathbf{Z}\) such that \(|x-k|_{p}\leqslant p^{-n}\) ; adding a suitable multiple of \(p^{n}\) to \(k\) , one can suppose that \(k\in[0,p^{n}-1]\) , and \(x\) is congruent to \(k\) modulo \(p^{n}\mathbf{Z}_{p}\) . Whence our assertion. This shows that \(\mathbf{Z}_{p}/p^{n}\mathbf{Z}_{p}\) is canonically isomorphic to \(\mathbf{Z}/p^{n}\mathbf{Z}\) . One sees, moreover, that \(\mathbf{Z}_{p}\) is precompact, hence compact since it is complete. Since \(\mathbf{Z}_{p}\) is an open subgroup of \(\mathbf{Q}_{p}\) , \(\mathbf{Q}_{p}\) is locally compact. The topology of \(\mathbf{Q}_{p}\) has a countable base (GT, IX, §2, No.~9, Cor. of Prop. 16). The additive group \(\mathbf{Q}_{p}\) may be identified with the inverse limit of the discrete groups \(\mathbf{Q}_{p}/p^{n}\mathbf{Z}_{p}\) .

There exists one and only one Haar measure \(\alpha\) on the additive group \(\mathbf{Q}_p\) such that \(\alpha(\mathbf{Z}_p) = 1\) ; it is called the normalized Haar measure on \(\mathbf{Q}_p\) . Since \(\mathbf{Z}_p\) is the union of \(p^n\) disjoint cosets of \(p^n\mathbf{Z}_p\) ( \(n\) an integer \(\geqslant 0\) ), one has \(\alpha(p^n\mathbf{Z}_p) = p^{-n}\) ; similarly \(\alpha(p^{-n}\mathbf{Z}_p) = p^n\) , so that, finally, \(\alpha(p^n\mathbf{Z}_p) = p^{-n}\) for every \(n \in \mathbf{Z}\) . By Prop. 8 b), \(\alpha\) is the only positive measure on \(\mathbf{Q}_p\) such that every coset of \(p^n\mathbf{Z}_p\) ( \(n\) an integer \(\geqslant 0\) ) has measure \(p^{-n}\) .

The restriction of \(\alpha\) to \(\mathbf{Z}_p\) is obviously a Haar measure on \(\mathbf{Z}_p\) .

